% https://tex.stackexchange.com/questions/286653/illustrating-stokes-theorem-with-asymptote
\documentclass[]{beamer}
\begin{document}
import graph3;
size(6cm,0);
settings.render = 4;
currentprojection=orthographic(4, 1, 1);

// COORDINATE SYSTEM
label("$O$", (0, 0, 0), NW);
Label Lx = Label("$x$", EndPoint, W);
Label Ly = Label("$y$", EndPoint, S);
Label Lz = Label("$z$", EndPoint, W);
draw(Lx, O--2X, Arrow3(emissive(black)));
draw(Ly, O--2Y, Arrow3(emissive(black)));
draw(Lz, O--2Z, Arrow3(emissive(black)));

// SURFACE
triple S(pair uv) {
  real x = cos(uv.x)*sin(uv.y);
  real y = sin(uv.x)*sin(uv.y);
  real z =           cos(uv.y);
  return (x-(1-z**2-y**2), y-0.3*sin(z*pi), z);
}

// BOUNDARY
triple dS(real u) {
  return S((u, pi/2));
}

// TANGENT VECTORS
real e = 0.001; // *this... is... infinitesimal!*
triple U(pair t) { return (S((e,0)+t)-S(t))/e; }
triple V(pair t) { return (S((0,e)+t)-S(t))/e; }

// NORMAL VECTOR
triple N(pair uv) {
  return cross(U(uv),V(uv))/length(cross(U(uv),V(uv)));
}

// GRAPHICAL ELEMENTS

surface s = surface(S, (0, 0), (2pi, pi/2), 64, 64, Spline);
path3 ds = graph(dS, 0, 2pi);

pair t = (-pi/8, pi/8);
path3 n = S(t) -- shift(-N(t))*S(t);

triple adj = (0,1.4,0.7);
s = shift(adj)*s;
ds = shift(adj)*ds;
n = shift(adj)*n;

label("$\Sigma$", shift(0,0,0.1)*shift(adj)*S((pi/2,pi/5)), NE);
Label LdS = Label("$\partial\Sigma$", Relative(0.95), S);
Label Ln = Label("$\mathbf n$", EndPoint, NW);

draw(s, surfacepen=material(diffusepen=gray(0.6),emissivepen=gray(0.3),specularpen=gray(0.1)));
draw(LdS, ds, Arrow3(emissive(black)));
draw(shift(4*unit(currentprojection.camera))*ds, dashed);
draw(Ln, n, Arrow3(emissive(black)));

\end{document}
